\chapter{Βελτιστοποίηση Πολυωνυμικού Πολλαπλασιασμού}
\label{ch:polymul-optimization}

Στο προηγούμενο κεφάλαιο τεκμηριώθηκε ότι ο πολυωνυμικός πολλαπλασιασμός
καταναλώνει $\approx\pct{97}$ του συνολικού \instret στο reference scalar
HQC-128, καθιστώντας τον μοναδικό κυρίαρχο bottleneck.  Το παρόν κεφάλαιο
παρουσιάζει τη συστηματική βελτιστοποίηση αυτού του πυρήνα σε τρία deployment
tiers αυξανόμενης αρχιτεκτονικής πολυπλοκότητας — από scalar Karatsuba με
\Zbc, μέσω διανυσματοποίησης recomposition με \RVV, έως πλήρη αξιοποίηση
\Zvbc στα leaf kernels.

%% ============================================================================
\section{Σχεδιαστικό Πλαίσιο: Hardware Tiers και Implementation Paths}
\label{sec:polymul-tiers}

Η modular φιλοσοφία του RISC-V \gls{isa} — κάθε επέκταση είναι προαιρετική
και ανεξάρτητα τυποποιημένη — δημιουργεί ένα φάσμα hardware capabilities στο
οποίο κάθε πλατφόρμα υλοποιεί διαφορετικό υποσύνολο επεκτάσεων.  Αντί μίας
μονολιθικής υλοποίησης που απαιτεί τη μέγιστη ISA, η παρούσα εργασία
στρωματοποιεί τη βελτιστοποίηση σε τρία \defterm{deployment tiers}, καθένα
από τα οποία αξιοποιεί αποκλειστικά ratified extensions και παράγει μετρήσιμη
επιτάχυνση σε σχέση με το reference baseline.

Η στρωματοποίηση αυτή εξυπηρετεί τρεις σκοπούς.  Πρώτον, εξασφαλίζει
portability: κάθε tier εκτελείται σε οποιονδήποτε επεξεργαστή υλοποιεί τις
αντίστοιχες επεκτάσεις, χωρίς custom silicon ή μη-τυποποιημένες εντολές.
Δεύτερον, επιτρέπει αξιολόγηση της οριακής συνεισφοράς κάθε επέκτασης:
η διαφορά μεταξύ Tier~1 και Tier~2 απομονώνει το κέρδος της διανυσματικής
μονάδας, ενώ η διαφορά μεταξύ Tier~2 και Tier~3 απομονώνει το κέρδος των
κρυπτογραφικών vector extensions.  Τρίτον, αντικατοπτρίζει τη ρεαλιστική
πορεία υιοθέτησης: σήμερα υπάρχει εμπορικό hardware για Tier~1 και Tier~2
(SpacemiT K1 X60, βλ.~\cref{ssec:bpi-f3}), ενώ το Tier~3 αντιστοιχεί στο
προφίλ RVA23 που βρίσκεται σε πορεία υιοθέτησης.

Ο \cref{tab:polymul-tiers} συνοψίζει τα τρία tiers, και το
\cref{fig:polymul-tiers} απεικονίζει τη σχέση υπερσυνόλου μεταξύ τους.

\begin{table}[htbp]
  \centering
  \caption{Deployment tiers πολυωνυμικού πολλαπλασιασμού.  Κάθε tier
    αξιοποιεί αποκλειστικά ratified RISC-V extensions.  Η στήλη
    \q{Champion} αναφέρει τη βελτιστοποιημένη υλοποίηση που αντιστοιχεί
    σε κάθε tier.}
  \label{tab:polymul-tiers}
  \begin{tabular}{clll}
    \toprule
    \textbf{Tier} & \textbf{ISA Profile} & \textbf{Κρίσιμες Εντολές} & \textbf{Champion} \\
    \midrule
    1 & \code{rv64gc\_zbc\_zbb\_zbs}
      & \clmul, \clmulh, \rori, \insn{andn}
      & \code{karatsuba\_zbc} \\
    2 & \code{rv64gcv\_zbc\_zbb\_zba}
      & + \insn{vle64}, \insn{vxor.vv}, \insn{vsetvl}
      & \code{karatsuba\_rvv} \\
    3 & \code{rv64gcv\_zbc\_zbb\_zba\_zvbc\_zvbb}
      & + \vclmul, \vclmulh, \vror
      & \code{karatsuba\_zvbc} \\
    \bottomrule
  \end{tabular}
\end{table}

\begin{figure}[htbp]
  \centering
  % Hardware deployment tiers — layered diagram
% Used in §6.1 (sec:polymul-tiers)

\begin{tikzpicture}[
  tier/.style={
    draw, rounded corners=4pt, minimum width=11cm, minimum height=1.6cm,
    align=center, font=\sffamily, thick
  },
  extlabel/.style={
    font=\scriptsize\ttfamily, anchor=west
  },
  arrow/.style={
    ->, >=stealth, thick, gray!60
  },
  brace/.style={
    decorate, decoration={brace, amplitude=6pt, mirror}, thick, gray!70
  }
]

% Tier 1 (base layer)
\node[tier, fill=blue!8, draw=blue!50] (t1) at (0, 0) {
  \textbf{Tier 1 — Scalar Deployment}\\[2pt]
  \footnotesize rv64gc + Zbc/Zbb/Zbs
};

% Tier 2 (middle layer)
\node[tier, fill=teal!8, draw=teal!50] (t2) at (0, 2.2) {
  \textbf{Tier 2 — Base-RVV Deployment}\\[2pt]
  \footnotesize rv64gcv + Zbc/Zbb/Zba
};

% Tier 3 (top layer)
\node[tier, fill=violet!8, draw=violet!50] (t3) at (0, 4.4) {
  \textbf{Tier 3 — Full Vector-Crypto Deployment}\\[2pt]
  \footnotesize rv64gcv + Zbc/Zbb/Zba/Zvbc/Zvbb
};

% Extension annotations (right side)
\node[extlabel] at (5.9, 0) {clmul, clmulh, rori, andn};
\node[extlabel] at (5.9, 2.2) {+ vle64, vse64, vxor, vsetvl};
\node[extlabel] at (5.9, 4.4) {+ vclmul.vv, vclmulh.vv, vror.vv};

% Variant labels (left side)
\node[font=\footnotesize\sffamily, anchor=east, blue!70!black] at (-5.9, 0)
  {\texttt{karatsuba\_zbc}};
\node[font=\footnotesize\sffamily, anchor=east, teal!70!black] at (-5.9, 2.2)
  {\texttt{karatsuba\_rvv}};
\node[font=\footnotesize\sffamily, anchor=east, violet!70!black] at (-5.9, 4.4)
  {\texttt{karatsuba\_zvbc}};

% Superset arrows
\draw[arrow] (t1.north) -- (t2.south)
  node[midway, right, font=\scriptsize, gray!80] {$+$ RVV 1.0};
\draw[arrow] (t2.north) -- (t3.south)
  node[midway, right, font=\scriptsize, gray!80] {$+$ Zvbc/Zvbb};

% Left brace with "ISA complexity increases" annotation
\draw[brace] (-6.2, -0.8) -- (-6.2, 5.2)
  node[midway, left=8pt, font=\scriptsize\sffamily, align=center, rotate=90]
  {Αρχιτεκτονική πολυπλοκότητα αυξάνεται};

\end{tikzpicture}

  \caption{Στρωματοποίηση deployment tiers ως σχέση υπερσυνόλου.  Κάθε
    ανώτερο tier περιλαμβάνει όλες τις επεκτάσεις του κατώτερου και
    προσθέτει νέες δυνατότητες.  Η αρχιτεκτονική πολυπλοκότητα αυξάνεται
    προς τα πάνω, ενώ η βάση (Tier~1) εξασφαλίζει τη μέγιστη portability.}
  \label{fig:polymul-tiers}
\end{figure}

Κοινός αλγοριθμικός πυρήνας και στα τρία tiers είναι ο Karatsuba με στατικά
παραγόμενο δέντρο αναδρομής (AP-2 generator, βλ.~\cref{ssec:ap2-generator}).
Η διαφοροποίηση εντοπίζεται σε δύο σημεία: (α)~τον τρόπο υλοποίησης του
\defterm{leaf kernel} — η βασική πράξη πολλαπλασιασμού στα φύλλα του δέντρου —
και (β)~τον τρόπο εκτέλεσης των \defterm{recomposition loops} — οι γραμμικού
κόστους βρόχοι XOR που συνθέτουν τα μερικά γινόμενα σε κάθε κόμβο αναδρομής.
Ο \cref{tab:tier-differentiation} αποσαφηνίζει αυτή τη διαφοροποίηση.

\begin{table}[htbp]
  \centering
  \caption{Διαφοροποίηση υλοποίησης μεταξύ tiers.  Ο αλγοριθμικός σκελετός
    (Karatsuba AP-2) είναι κοινός — μόνο τα leaf kernels και οι recomposition
    loops διαφέρουν.}
  \label{tab:tier-differentiation}
  \begin{tabular}{lll}
    \toprule
    \textbf{Tier} & \textbf{Leaf Kernel} & \textbf{Recomposition Loops} \\
    \midrule
    1 & Scalar \clmul/\clmulh (Zbc)
      & Scalar unrolled XOR \\
    2 & Scalar \clmul/\clmulh (Zbc)
      & Inlined RVV (\insn{vxor.vv}, tuned LMUL) \\
    3 & Vector \vclmul/\vclmulh (Zvbc)
      & Inlined RVV (\insn{vxor.vv}, tuned LMUL) \\
    \bottomrule
  \end{tabular}
\end{table}

Η αρχιτεκτονική αυτή παρατήρηση εξηγεί εκ των προτέρων τη σχετική
αποτελεσματικότητα κάθε tier.  Στο Tier~2, η διανυσματοποίηση επιταχύνει
μόνο τους γραμμικού κόστους $\BigO{n}$ recomposition loops, ενώ ο κυρίαρχος
τετραγωνικού κόστους $\BigO{n^2}$ leaf παραμένει scalar — το κέρδος είναι
πραγματικό αλλά δομικά περιορισμένο.  Στο Tier~3, η διανυσματοποίηση
εισχωρεί στον ίδιο τον leaf, μειώνοντας τον αριθμό εντολών ανά
\num{64}-bit carry-less γινόμενο και αποδίδοντας ουσιαστικά μεγαλύτερη
επιτάχυνση.

Οι επόμενες ενότητες παρουσιάζουν κάθε tier αναλυτικά: η
\cref{sec:polymul-scalar} αναπτύσσει το scalar path (Tier~1), η
\cref{sec:polymul-base-rvv} τη διανυσματοποίηση recomposition (Tier~2),
και η \cref{sec:polymul-zvbc} τα vector-crypto leaf kernels (Tier~3).

%% ----------------------------------------------------------------------------
\subsection{Tier 1 — Scalar Deployment (rv64gc + Zbc/Zbb)}
\label{ssec:tier1}

Το Tier~1 στοχεύει πλατφόρμες χωρίς διανυσματική μονάδα — ενσωματωμένους
επεξεργαστές και IoT κόμβους που υλοποιούν μόνο τις scalar bit-manipulation
extensions.  Η κρίσιμη αρχιτεκτονική δυνατότητα είναι η εντολή \clmul{} της
\Zbc: ένα ζεύγος \clmul/\clmulh υπολογίζει το πλήρες \num{128}-bit carry-less
γινόμενο δύο \num{64}-bit πολυωνύμων σε δύο εντολές μηχανής, αντικαθιστώντας
τον bit-serial βρόχο \num{64} conditional shift-XOR iterations του reference
implementation.

Η βελτιστοποίηση Tier~1 συνδυάζει τρεις τεχνικές:
%
\begin{enumerate}
  \item \textbf{Word-level Karatsuba leaf:} Αντικατάσταση του bit-serial
    schoolbook με word-granularity schoolbook βασισμένο σε \clmul/\clmulh
    (\cref{ssec:bit-to-word}).

  \item \textbf{Assign-on-first-write discipline:} Εξάλειψη του \code{memset}
    μέσω πειθαρχίας πρώτης εγγραφής — κάθε output word λαμβάνει direct
    assignment στην πρώτη πρόσβαση (\cref{ssec:memset-elimination}).

  \item \textbf{Στατική παραγωγή δέντρου αναδρομής (AP-2):} Compile-time
    generation ολόκληρου του Karatsuba recursion tree, εξαλείφοντας runtime
    dispatch, stack spilling, και ABI overhead
    (\cref{ssec:ap2-generator}).
\end{enumerate}

\noindent
Η ISA string \code{rv64gc\_zbc\_zbb\_zbs} μεταγλωττίζεται με
\code{-fno-tree-vectorize}, εξασφαλίζοντας ότι ο compiler δεν εκπέμπει
RVV εντολές — η μέτρηση αντικατοπτρίζει αποκλειστικά scalar εκτέλεση.
Η \Zbb παρέχει \rori{} (native rotate) και \insn{andn} (AND-NOT) που
αξιοποιούνται αυτόματα μέσω compiler auto-lowering χωρίς αλλαγή πηγαίου
κώδικα.

%% ----------------------------------------------------------------------------
\subsection{Tier 2 — Base-RVV Deployment (rv64gcv χωρίς Zvbc/Zvbb)}
\label{ssec:tier2}

Το Tier~2 στοχεύει application-class πλατφόρμες με διανυσματική μονάδα
RVV~1.0 αλλά χωρίς τις κρυπτογραφικές vector extensions.  Η ISA string
\code{rv64gcv\_zbc\_zbb\_zba} προσθέτει τη δυνατότητα εκτέλεσης
data-parallel πράξεων σε ευρείς vector registers (\VLEN bits ανά register,
τυπικά \num{128}–\num{2048}).

Η αρχιτεκτονική παρατήρηση-κλειδί είναι ότι η \Zvbc δεν είναι διαθέσιμη
σε αυτό το tier: η εντολή \vclmul{} απουσιάζει, επομένως ο leaf kernel
παραμένει scalar (\clmul/\clmulh, ταυτόσημος με Tier~1).  Η διανυσματική
μονάδα αξιοποιείται αποκλειστικά στους recomposition loops — τους γραμμικού
κόστους βρόχους XOR, copy, και modular reduction που εκτελούνται σε κάθε
κόμβο του δέντρου αναδρομής.

Η βελτιστοποίηση Tier~2 εισάγει τρεις τεχνικές πέραν του Tier~1:
%
\begin{enumerate}
  \item \textbf{Direct-write discipline:} Τα child multiplications γράφουν
    απευθείας στο output buffer, εξαλείφοντας δύο πλήρεις RVV copy loops
    ανά κόμβο αναδρομής (\cref{ssec:direct-write}).

  \item \textbf{Overlap-save combine:} Επίλυση του aliasing hazard που
    εισάγει η direct-write μέσω αποθήκευσης της overlap region σε scratch
    buffer (\cref{ssec:overlap-save}).

  \item \textbf{Inlined RVV operations με tuned LMUL:} Αντικατάσταση
    shared helper functions με inline RVV intrinsic loops, με \LMUL
    ρυθμισμένο ανά operation βάσει register pressure
    (\cref{ssec:inlined-rvv}).
\end{enumerate}

\noindent
Η δομική αυτή αρχιτεκτονική — scalar leaf + vector recomposition — αποτελεί
\defterm{υβριδική προσέγγιση}: η αριθμητική πράξη (carry-less multiply)
εκτελείται scalar, ενώ η data-movement πράξη (XOR, copy, reduce)
εκτελείται vector.  Η προσέγγιση αυτή αποδίδει μετρήσιμη βελτίωση
(\cref{sec:polymul-base-rvv}), αλλά το κέρδος περιορίζεται δομικά από
τον scalar leaf — ο τετραγωνικός $\BigO{n^2}$ leaf κυριαρχεί έναντι
των γραμμικών $\BigO{n}$ recomposition loops.

%% ----------------------------------------------------------------------------
\subsection{Tier 3 — Full Vector-Crypto Deployment (rv64gcv + Zvbc/Zvbb)}
\label{ssec:tier3}

Το Tier~3 αντιπροσωπεύει την πλήρη αξιοποίηση κρυπτογραφικών vector
extensions.  Η ISA string \code{rv64gcv\_zbc\_zbb\_zba\_zvbc\_zvbb} προσθέτει
τις εντολές \vclmul{} και \vclmulh{} (Zvbc) — vector carry-less multiply
που υπολογίζει πολλαπλά \num{64}$\times$\num{64}-bit carry-less γινόμενα
παράλληλα σε ένα vector register group.

Η θεμελιώδης αρχιτεκτονική αλλαγή σε σχέση με το Tier~2 είναι η
\defterm{διανυσματοποίηση του leaf kernel}: ο scalar \clmul/\clmulh
αντικαθίσταται από τον \code{vclmul\_leaf}, ο οποίος εκτελεί ολόκληρο
τον schoolbook πολλαπλασιασμό μεγέθους $n=16$ words μέσω vector
instructions.  Σε \VLEN=\num{256} με \LMUL=m4, κάθε vector register group
περιέχει \num{16} elements — ακριβώς το μέγεθος του Karatsuba threshold —
επομένως ο εσωτερικός βρόχος συμπτύσσεται σε μία μόνο vector iteration
ανά εξωτερικό index.

Η \Zvbb παρέχει επιπλέον τις εντολές \vror{} (vector rotate) και
\insn{vandn.vv} (vector AND-NOT), οι οποίες αξιοποιούνται στον Keccak
πυρήνα (\cref{ch:kernel-optimization}) αλλά δεν επηρεάζουν άμεσα τον
πολυωνυμικό πολλαπλασιασμό.

Η σημασία του Tier~3 εκτείνεται πέραν της ποσοτικής επιτάχυνσης.
Η \cref{sec:isa-ablation} τεκμηρίωσε
ότι ο compiler δεν συνθέτει αυτόματα \vclmul{} από generic C loops —
η διαφορά μεταξύ \code{ref\_vector} (compiler ceiling) και
\code{karatsuba\_zvbc} (χειρόγραφα intrinsics) αποτελεί αποκλειστικά
αποτέλεσμα ρητής αλγοριθμικής παρέμβασης.  Η παρατήρηση αυτή ενισχύει
τη θέση ότι οι κρυπτογραφικές vector extensions αποδίδουν μόνο μέσω
εξειδικευμένων υλοποιήσεων — η ύπαρξη hardware δεν αρκεί χωρίς
αντίστοιχο software.

\bigskip
\noindent
Η ορθότητα κάθε tier επαληθεύεται ανεξάρτητα μέσω του πλήρους KAT suite
(\num{300} vectors, τρία security levels) — κανένα tier δεν εισάγει
αλγοριθμική απόκλιση από τη reference υλοποίηση.  Η μεταγλώττιση κάθε
variant με αυστηρά ελεγχόμενη \code{-march} string
(\cf \cref{tab:march-per-variant}) εξασφαλίζει ότι ο assembler απορρίπτει
κάθε εντολή εκτός tier — ένα \vclmul{} σε Tier~1 target προκαλεί σφάλμα
μεταγλώττισης, αποτρέποντας ακούσια εξάρτηση από μη-διαθέσιμο hardware.

%% ============================================================================
\section{Scalar Path: Karatsuba με \Zbc}
\label{sec:polymul-scalar}

\subsection{Μετάβαση από Bit-Serial σε Word-Level Πολλαπλασιασμό}
\label{ssec:bit-to-word}

\subsection{Assign-on-First-Write: Εξάλειψη memset}
\label{ssec:memset-elimination}

\subsection{Δομή Αναδρομής και AP-2 Generator}
\label{ssec:ap2-generator}

\subsection{Απορριφθείσες Εναλλακτικές (Toom-Cook, FAFFT)}
\label{sec:rejected-alternatives}

%% ============================================================================
\section{Base-RVV Path: Διανυσματοποίηση Recomposition}
\label{sec:polymul-base-rvv}

\subsection{Direct-Write Discipline: Εξάλειψη Αντιγραφών z0/z2}
\label{ssec:direct-write}

\subsection{Overlap-Save Combine Strategy}
\label{ssec:overlap-save}

\subsection{Inlined RVV Operations με Tuned LMUL}
\label{ssec:inlined-rvv}

%% ============================================================================
\section{Vector-Crypto Path: \Zvbc Leaf Kernels}
\label{sec:polymul-zvbc}

%% ============================================================================
\section{Code Size vs Performance Trade-off}
\label{sec:polymul-codesize}
